% 同一十封邮件，9、10被分错；两行是相同样本的两种完整排序。
\begin{tikzpicture}[figure picture,foundation picture]
 \path[use as bounding box] (0,0) rectangle (168,65);
 \foreach \sy/\name/\risk in {0/甲/0,-27/乙/25\%}{
  \begin{scope}[yshift=\sy mm]
   \node[anchor=west] at (0,59) {\name};
   \draw[foundation axis] (22,59)--(94,59);
   \node[foundation edge label] at (57,59) {$u(\bm x)$};
   \path[fill=FigureBlueFill,rounded corners=1mm] (0,40) rectangle (74,55);
   \path[fill=FigurePaper,rounded corners=1mm] (75,40) rectangle (94,55);
   \draw[draw=FigureStroke,dashed,line width=.8pt] (74.5,38)--(74.5,56);
   \node[anchor=west] at (76,58) {$t$};
   \node at (37,35) {$g_t=1$};
   \node at (85,35) {$g_t=0$};
   \node[anchor=west] at (103,53) {$\hat c(t)=\dfrac{8}{10}=80\%$};
   \node[anchor=west] at (103,42)
    {$\hat R_{\mathrm{sel}}(t)=\dfrac{\ifnum\sy=0 0\else 2\fi}{8}=\risk$};
  \end{scope}
 }
 \foreach \i/\id/\wrong in {0/1/0,1/2/0,2/3/0,3/4/0,4/5/0,5/6/0,6/7/0,7/8/0,8/9/1,9/10/1}{
  \begin{scope}[xshift={(\i*9.3+1)*1mm}]
   \draw[draw=FigureStroke,fill=white,rounded corners=.5mm,line width=1pt]
    (0,47) rectangle (7.6,54);
   \node at (3.8,50.5) {\id};
   \ifnum\wrong=1 \node[text=FigureRed] at (3.8,43) {$\times$};
   \else \node[text=FigureBlue] at (3.8,43) {$\checkmark$}; \fi
  \end{scope}
 }
 \foreach \i/\id/\wrong in {0/9/1,1/10/1,2/3/0,3/4/0,4/5/0,5/6/0,6/7/0,7/8/0,8/1/0,9/2/0}{
  \begin{scope}[xshift={(\i*9.3+1)*1mm},yshift=-27mm]
   \draw[draw=FigureStroke,fill=white,rounded corners=.5mm,line width=1pt]
    (0,47) rectangle (7.6,54);
   \node at (3.8,50.5) {\id};
   \ifnum\wrong=1 \node[text=FigureRed] at (3.8,43) {$\times$};
   \else \node[text=FigureBlue] at (3.8,43) {$\checkmark$}; \fi
  \end{scope}
 }
\end{tikzpicture}
